\chapter{Core Source Code}
\label{ch:appendix_code}

\section{Outlier Filtering Logic}
The python source code for univariate outlier screening inside \textit{OutlierEngine} is shown below:

\begin{lstlisting}[language=python, caption=Univariate Outlier Detection Code]
def detect(self, ratio_matrix: pd.DataFrame) -> pd.DataFrame:
    records = []
    ratio_cols = [c for c in ratio_matrix.columns if c not in EXCLUDE_COLS]

    for ratio in ratio_cols:
        col = ratio_matrix[ratio].dropna()
        if len(col) < self.min_peers:
            continue

        q1, q3 = col.quantile(0.25), col.quantile(0.75)
        iqr = q3 - q1
        lower = q1 - self.iqr_multiplier * iqr
        upper = q3 + self.iqr_multiplier * iqr
        median = col.median()
        mean = col.mean()

        # Robust Z-Scores
        mad = (col - median).abs().median()
        if mad == 0:
            std = col.std()
            zscore_series = (col - mean) / std if std != 0 else pd.Series(0.0, index=col.index)
        else:
            zscore_series = 0.6745 * (col - median) / mad

        for ticker in ratio_matrix.index:
            val = ratio_matrix.at[ticker, ratio]
            if pd.isna(val):
                continue

            iqr_flag = bool(val < lower or val > upper)
            z = float(zscore_series.get(ticker, 0.0))
            zscore_flag = bool(abs(z) > self.zscore_threshold)

            if iqr_flag or zscore_flag:
                dev_pct = ((val - median) / abs(median) * 100) if median != 0 else np.nan
                severity = self._severity(val, lower, upper, iqr, z)
                direction = self._direction(ratio, val, median)
                is_positive = self._is_positive_outlier(ratio, direction)

                records.append({
                    "ticker": ticker,
                    "ratio": ratio,
                    "value": round(val, 4),
                    "sector_median": round(median, 4),
                    "zscore": round(z, 3),
                    "iqr_flag": iqr_flag,
                    "zscore_flag": zscore_flag,
                    "deviation_pct": round(dev_pct, 2) if not pd.isna(dev_pct) else np.nan,
                    "severity_score": round(severity, 2),
                    "is_positive_outlier": is_positive,
                })

    return pd.DataFrame(records)
\end{lstlisting}

\chapter{Tally ERP Ingestion Formats}
\label{ch:appendix_tally_formats}
The required CSV structures for uploading private enterprise ledger groups and double-entry transaction journals in the \textit{TallySync} bridge are defined below:

\section{Ledger Master CSV Format}
The Ledger Master file maps local transaction accounts to their primary Tally accounting groups:

\begin{lstlisting}[language=python, caption=Ledger Master CSV Format]
CLIENT_ID,COMPANY_NAME,FINANCIAL_YEAR,NAME,PARENT,PRIMARYGROUP,OPENING,CLOSING,GUID,SYNC_TIME
10000_VSM,VIBRANT SUPERMARKET,2025-26,0% CST SALES,CST T/O,Sales Accounts,0.0,0.0,783aa38d-...,15-04-26 20:26
10000_VSM,VIBRANT SUPERMARKET,2025-26,0% Integrated Supply,Purchase Accounts,Purchase Accounts,0.0,-6000907.42,783aa38d-...,15-04-26 20:26
\end{lstlisting}

\section{Journal Voucher CSV Format}
The Journal Voucher file streams the actual credit and debit transaction journal voucher records:

\begin{lstlisting}[language=python, caption=Journal Voucher CSV Format]
CLIENT_ID,COMPANY_NAME,FINANCIAL_YEAR,DATE,PARTICULARS,VCHTYPE,VCHNO,DEBIT,CREDIT,VOUCHERKEY,NARRATION
10000_VSM,VIBRANT SUPERMARKET,2025-26,1-Apr-25,KISHOREKUMAR,GST Common Invoices,KIS-01/25-26,0.0,2000.0,783aa38d-...,Being Rent
10000_VSM,VIBRANT SUPERMARKET,2025-26,1-Apr-25,Rent,GST Common Invoices,KIS-01/25-26,-2000.0,0.0,783aa38d-...,Being Rent
\end{lstlisting}

